\documentclass[aspectratio=169, 11pt, xcolor={table}]{beamer}
% Shared style for the course slides, based on presentation_text/main.tex.
% Dark blue removed: every slide (cover included) uses the light background
% of the template's content slides, with teal + orange as accents.

\usepackage{amsmath}
\usepackage{amssymb}
\usepackage{booktabs}
\usepackage{colortbl}
\usepackage{tabularx}
\usepackage{tikz}
\usepackage[skins]{tcolorbox}
\usepackage{fontawesome5}
\usepackage{polyglossia}

\usetikzlibrary{positioning, calc, arrows.meta, shapes.geometric}

%------------------------------------------------
% Languages and fonts (compile from slides/ with XeLaTeX or tectonic)
%------------------------------------------------
\setdefaultlanguage[numerals=maghrib]{arabic}
\setotherlanguage{english}
\usefonttheme{serif}

\newfontfamily\arabicfont[
    Script=Arabic, Scale=0.98, Path=../fonts/, Extension=.ttf,
    UprightFont=*-Regular, BoldFont=*-Bold,
    ItalicFont=*-Italic, SlantedFont=*-Italic,
    BoldItalicFont=*-BoldItalic, BoldSlantedFont=*-BoldItalic]{Amiri}
\newfontfamily\arabicfontsf[
    Script=Arabic, Scale=0.98, Path=../fonts/, Extension=.ttf,
    UprightFont=*-Regular, BoldFont=*-Bold,
    ItalicFont=*-Italic, SlantedFont=*-Italic,
    BoldItalicFont=*-BoldItalic, BoldSlantedFont=*-BoldItalic]{Amiri}
% Heading font
\newfontfamily\khatt[Script=Arabic, Scale=1.05, Path=../fonts/,
    BoldFont=ae_AlBattar.ttf]{ae_AlBattar.ttf}

%------------------------------------------------
% Colors (template palette without the navy tones)
%------------------------------------------------
\definecolor{brandTeal}{HTML}{00B4D8}
\definecolor{brandOrange}{HTML}{FF7043}
\definecolor{brandLight}{HTML}{F8F9FA}
\definecolor{brandDark}{HTML}{263238}
\definecolor{brandMuted}{HTML}{6C757D}
\colorlet{tealDeep}{brandTeal!70!black}

\setbeamercolor{background canvas}{bg=brandLight}
\setbeamercolor{normal text}{fg=brandDark}
\setbeamercolor{structure}{fg=brandTeal}
\setbeamertemplate{navigation symbols}{}

% beamer.cls sets \raggedright globally; under RTL that centers short lines
% or pushes them left, so go back to justified text (flush right).
\rightskip=0pt\relax
% Lists call \raggedright too; make it a no-op for RTL slides.
\AtBeginDocument{\renewcommand{\raggedright}{\rightskip=0pt\relax}}

%------------------------------------------------
% Frame title, footer, lists
%------------------------------------------------
\setbeamertemplate{frametitle}{%
    \vspace{0.35cm}
    \noindent{\khatt\bfseries\Large\color{brandDark}\insertframetitle}
    \par
    \vskip-2pt
    \noindent{\color{brandTeal}\rule{\textwidth}{1.5pt}}
    \vspace{0.05cm}
}

% Footer: page number only
\setbeamertemplate{footline}{%
    \hbox to \paperwidth{\hfill
        {\color{brandMuted}\scriptsize\LR{\insertframenumber\,/\,\inserttotalframenumber}}%
        \hspace{0.6cm}}%
    \vspace{0.25cm}}

\setbeamertemplate{itemize item}{\color{brandTeal}\small\faChevronLeft}
\setbeamertemplate{itemize subitem}{\color{brandMuted}\small\faAngleLeft}
\setbeamertemplate{enumerate item}{\color{tealDeep}\bfseries\insertenumlabel)}

% Tables: every cell framed by a visible border
\newcolumntype{R}{>{\raggedleft\arraybackslash}X}
\newcolumntype{C}{>{\centering\arraybackslash}X}
\newcommand{\ras}[1]{\cellcolor{brandTeal!20}\textbf{#1}}
% Colored text inside a table cell (\leavevmode keeps it on the row baseline)
\newcommand{\cell}[2]{\leavevmode\textcolor{#1}{#2}}
\renewcommand{\arraystretch}{1.35}
\setlength{\arrayrulewidth}{0.6pt}
\arrayrulecolor{brandMuted!70}

%------------------------------------------------
% Boxes (the two boxes of the template's 2nd slide)
%------------------------------------------------
\tcbset{
    brandbox/.style={
        enhanced, arc=4pt, outer arc=4pt, boxrule=1.5pt,
        coltitle=white, coltext=brandDark, fonttitle=\khatt\bfseries,
        toptitle=0.8mm, bottomtitle=0.8mm,
        top=1.2mm, bottom=1.2mm, left=3mm, right=3mm,
        before skip=4pt, after skip=4pt}
}
% Teal box: definitions, objectives, summaries
\newtcolorbox{tbox}[2][]{brandbox, colback=brandTeal!10,
    colframe=brandTeal, title={#2}, #1}
% Orange box: warnings
\newtcolorbox{obox}[2][]{brandbox, colback=brandOrange!10,
    colframe=brandOrange, title={#2}, #1}
% Card (template's 3rd slide, in teal)
\newtcolorbox{card}[2][]{enhanced, colback=brandTeal!6, colframe=brandTeal,
    arc=6pt, boxrule=1.2pt, halign=center, halign title=center,
    coltitle=white, coltext=brandDark, fonttitle=\khatt\bfseries,
    top=1.5mm, bottom=1.5mm, left=2mm, right=2mm, title={#2}, #1}
% Prompt text written to the tool
\newtcolorbox{amr}[1][brandTeal]{enhanced, arc=1mm, boxrule=0pt,
    colback=white, overlay={\fill[#1] (frame.north east) rectangle ([xshift=-2.5pt]frame.south east);},
    left=6pt, right=8pt, top=2pt, bottom=2pt,
    before skip=3pt, after skip=3pt}

%------------------------------------------------
% TikZ (Arabic right-to-left inside nodes)
%------------------------------------------------
\newcommand{\nodecolor}{}
\tikzset{
    every node/.append style={
        execute at begin node={\nodecolor\everypar{\beginR}\beginR},
        execute at end node=\endR},
    lawn/.code={\def\nodecolor{\color{#1}}},
    khatwa/.style={draw=#1, fill=#1!12, very thick, rounded corners=5pt,
        minimum height=0.95cm, text width=1.9cm, align=center,
        font=\khatt\bfseries\small},
    khatwa/.default=brandTeal,
    sahm/.style={-{Stealth[length=2.5mm]}, thick, color=brandMuted},
    rang/.style={draw=#1, fill=#1!12},
    raqm/.style={circle, fill=#1, lawn=white, inner sep=1.5pt,
        font=\small\bfseries},
}

% Horizontal chain of steps, right to left:
% \flow[width]{{a}, {b\\c}, {d}}
\newcommand{\flow}[2][1.9cm]{%
    \begin{tikzpicture}[node distance=5mm]
        \foreach \t [count=\i] in {#2} {
            \ifnum\i=1
                \node[khatwa, text width=#1] (f\i) {\t};
            \else
                \pgfmathtruncatemacro\p{\i-1}
                \node[khatwa, text width=#1, left=of f\p] (f\i) {\t};
                \draw[sahm] (f\p) -- (f\i);
            \fi
        }
    \end{tikzpicture}}

%------------------------------------------------
% Cover (template's 1st slide on the light background)
%------------------------------------------------
\newcommand{\cover}[3]{%
\begin{frame}[plain]
    \begin{tikzpicture}[remember picture, overlay]
        \draw[brandTeal, line width=1.5pt, opacity=0.3]
            (current page.south east) ++(-1,1) circle (4cm);
        \draw[brandTeal, line width=0.8pt, opacity=0.25]
            (current page.north west) ++(2,2) circle (6cm);
        \draw[brandOrange, line width=0.5pt, opacity=0.25]
            (current page.north west) ++(2,2) circle (7cm);
        \node[anchor=north east] at ($(current page.north east)+(-1.0,-0.8)$) {
            \begin{tikzpicture}[scale=0.35]
                \fill[brandTeal] (0,0) -- (0.8,1.2) -- (1.6,0) -- cycle;
                \fill[brandOrange] (0.8,1.2) -- (1.6,2.4) -- (2.4,1.2) -- cycle;
            \end{tikzpicture}
        };
    \end{tikzpicture}

    \vspace{-0.05cm}
    {\hspace*{1.1cm}\khatt\bfseries\color{brandDark}
        المدرسة العليا للأساتذة ببوزريعة ــ ملحقة بومرداس\par}

    \vspace{1.2cm}
    \begin{minipage}{0.85\textwidth}
        {\color{tealDeep}\khatt\bfseries #1} \\[0.3cm]
        {\raggedleft\color{brandDark}\Huge\khatt\bfseries #2\par} \vspace{0.3cm}
        {\color{brandMuted}\large #3}
    \end{minipage}

    \vspace{1.2cm}
    \begin{minipage}{0.85\textwidth}
        \color{brandMuted}\small
        \begin{tabular}{@{}ll}
            \textbf{الأستاذ:} & د. بن عبد الرزاق يوسف \\
            \textbf{الفئة:}   & طلبة قسم اللغة العربية \\
        \end{tabular}
    \end{minipage}
\end{frame}}

\title{الذكاء الاصطناعي متعدد الأنماط}

\begin{document}

\cover{درس الفصل السادس}
{الذكاء الاصطناعي متعدد الأنماط}
{نص ــ صورة ــ صوت ــ فيديو}

\begin{frame}{أهداف الدرس}
    \begin{tbox}{\faBullseye\quad أهداف الدرس}
        في نهاية هذا الدرس، تستطيع أن:
        \begin{itemize}
            \item تشرح معنى «النمط» و«متعدد الأنماط».
            \item تستعمل أداة لفهم صورة أو مخطط، ثم تتحقق من إجابتها.
            \item تحوّل نصًا إلى مخطط مفاهيمي بمساعدة الأداة.
            \item تراجع المخطط وتصحح أخطاءه.
        \end{itemize}
    \end{tbox}
\end{frame}

\begin{frame}{ما معنى «متعدد الأنماط»؟}
    نتواصل نحن البشر بطرق كثيرة: نكتب، ونتكلم، ونرسم، ونشير.
    كل طريقة من هذه الطرق نسميها \textbf{نمطًا}.
    \begin{tbox}{\faLightbulb\quad تعريف}
        \textbf{النمط:} شكل المعلومة: نص، أو صورة، أو صوت، أو فيديو.\\[3pt]
        \textbf{الذكاء الاصطناعي متعدد الأنماط:} أداة تفهم أكثر من نمط،
        وتستطيع أن تنتج أكثر من نمط.
    \end{tbox}
\end{frame}

\begin{frame}{من المدخلات إلى المخرجات}
    \begin{center}
    \begin{tikzpicture}[
        namat/.style={draw=brandTeal, fill=brandTeal!12, thick, rounded corners=4pt,
            minimum width=2.6cm, minimum height=0.7cm, font=\khatt\bfseries}]
        \node[font=\khatt\bfseries, lawn=brandMuted] at (4.5,1.9) {المدخلات};
        \node[namat] (i1) at (4.5,1.1)   {نص};
        \node[namat] (i2) at (4.5,0.25)  {صورة};
        \node[namat] (i3) at (4.5,-0.6)  {صوت};
        \node[namat] (i4) at (4.5,-1.45) {فيديو};
        \node[draw=brandOrange, fill=brandOrange!12, very thick, rounded corners=6pt,
            minimum width=2.6cm, minimum height=3.2cm, align=center,
            font=\khatt\bfseries] (m) at (0,-0.17) {أداة\\متعددة\\الأنماط};
        \node[font=\khatt\bfseries, lawn=brandMuted] at (-4.5,1.9) {المخرجات};
        \node[namat] (o1) at (-4.5,1.1)   {شرح مكتوب};
        \node[namat] (o2) at (-4.5,0.25)  {مخطط};
        \node[namat] (o3) at (-4.5,-0.6)  {صورة جديدة};
        \node[namat] (o4) at (-4.5,-1.45) {قراءة صوتية};
        \foreach \i in {1,...,4} {
            \draw[sahm] (i\i.west) -- (m.east);
            \draw[sahm] (m.west) -- (o\i.east);
        }
    \end{tikzpicture}
    \end{center}
\end{frame}

\begin{frame}{أمثلة لأستاذ اللغة العربية}
    \begin{tbox}{\faChalkboardTeacher\quad في قسم اللغة العربية}
        \begin{itemize}
            \item \textbf{صورة $\leftarrow$ نص:} يصوّر صفحة مخطوطة، ويطلب كتابتها نصًا رقميًا.
            \item \textbf{نص $\leftarrow$ مخطط:} يحوّل درس «الجملة الاسمية» إلى خريطة ذهنية.
            \item \textbf{نص $\leftarrow$ صوت:} يستمع إلى قراءة لقصيدة قبل تقديمها للتلاميذ.
        \end{itemize}
    \end{tbox}
\end{frame}

\begin{frame}{فهم صورة باستعمال الذكاء الاصطناعي}
    نرفع إلى الأداة صورة فيها معلومات: مخطط بياني، أو جدول، أو رسم توضيحي،
    ثم نطلب منها شرح ما تراه.
    \begin{amr}[brandTeal]
        حلّل هذه الصورة، واشرح لي المعلومات الموجودة فيها بطريقة بسيطة.
    \end{amr}
    \textbf{\color{tealDeep}خطوات العمل:}
    \begin{center}
        \flow[2.2cm]{{اختر صورة\\واضحة}, {ارفعها\\مع أمر دقيق}, {اقرأ\\الشرح}, {قارن\\بالصورة}}
    \end{center}
\end{frame}

\begin{frame}{أين تخطئ الأداة غالبًا؟}
    \begin{obox}{\faExclamationTriangle\quad أخطاء متكررة}
        \begin{itemize}
            \item \textbf{الأرقام:} قد تقرأ رقمًا خطأ، أو تخلط بين عمودين.
            \item \textbf{العلاقات:} قد تفهم العلاقة بين عنصرين بشكل خاطئ.
            \item \textbf{الخط العربي:} قد تخطئ في قراءة خط اليد أو الحروف غير الواضحة.
        \end{itemize}
    \end{obox}
    لذلك نقارن دائمًا إجابة الأداة \textbf{بالصورة الأصلية}.
\end{frame}

\begin{frame}{تحويل النص إلى مخطط مفاهيمي}
    \begin{tbox}{\faLightbulb\quad تعريف}
        \textbf{المخطط المفاهيمي:} رسم يعرض المفاهيم الأساسية لدرس ما،
        ويربط بينها بأسهم توضح العلاقات.
    \end{tbox}
    \textbf{\color{tealDeep}الأمر:}
    \begin{amr}[brandTeal]
        حوّل النص الآتي إلى مخطط مفاهيمي بسيط، يوضح العلاقة بين المفاهيم الرئيسية.
    \end{amr}
\end{frame}

\begin{frame}{مثال: مخطط لدرس «الجملة الاسمية»}
    \begin{center}
    \begin{tikzpicture}[
        mafhum/.style={thick, rounded corners=4pt,
            minimum height=0.8cm, text width=2.4cm, align=center,
            font=\khatt\bfseries},
        alaqa/.style={-{Stealth[length=2.5mm]}, thick, brandMuted}]
        \node[mafhum, rang=brandOrange, text width=3.2cm] (j) at (0,0) {الجملة الاسمية};
        \node[mafhum, rang=brandTeal] (m) at (3,-1.7)  {المبتدأ};
        \node[mafhum, rang=brandTeal] (k) at (-3,-1.7) {الخبر};
        \node[mafhum, rang=brandMuted] (m1) at (3,-3.4) {اسم مرفوع};
        \node[mafhum, rang=brandMuted] (k1) at (-1.6,-3.4) {مفرد};
        \node[mafhum, rang=brandMuted] (k2) at (-4.6,-3.4) {جملة أو\\شبه جملة};
        \draw[alaqa] (j) -- node[above, sloped, font=\small] {تتكون من} (m);
        \draw[alaqa] (j) -- node[above, sloped, font=\small] {تتكون من} (k);
        \draw[alaqa] (m) -- node[right, font=\small] {هو} (m1);
        \draw[alaqa] (k) -- node[right, font=\small] {يكون} (k1);
        \draw[alaqa] (k) -- (k2);
    \end{tikzpicture}
    \end{center}
\end{frame}

\begin{frame}{مراجعة المخطط}
    بعد أن تقترح الأداة المخطط، نسأل أنفسنا:
    \begin{tbox}{\faQuestionCircle\quad أسئلة المراجعة}
        \begin{enumerate}
            \item هل يحتوي المخطط على \textbf{كل} الأفكار الأساسية في النص؟
            \item هل توجد علاقة \textbf{خاطئة} بين بعض المفاهيم؟
            \item هل أضافت الأداة مفهومًا \textbf{غير موجود} في النص؟
            \item ماذا أضيف أو أحذف ليصبح المخطط أوضح لتلاميذي؟
        \end{enumerate}
    \end{tbox}
\end{frame}

\begin{frame}{خلاصة الدرس}
    \begin{tbox}{\faCheckCircle\quad خلاصة الدرس}
        \begin{itemize}
            \item \textbf{النمط} هو شكل المعلومة: نص، صورة، صوت، فيديو.
            \item الأداة متعددة الأنماط تفهم أكثر من نمط وتنتج أكثر من نمط.
            \item الأداة قد تخطئ في \textbf{الأرقام والعلاقات}، فنقارن بالأصل دائمًا.
            \item المخطط المفاهيمي يحتاج إلى \textbf{مراجعة} الأستاذ قبل استعماله في القسم.
        \end{itemize}
    \end{tbox}
\end{frame}

\end{document}
